\section{Web Accessibility Research}
\label{research:WebAccessibility}
\textbf{Source:} \cite{UnderstandingWCAG222024}

Making the internet as a whole accessible to everyone should be a fundamental driving goal behind every web developer and designers creations. Available to all the internet hosts all manor of resource which can be hugely beneficial to a users life (with my artifact hopefully joining this group). 

Currently I know that common access requirements can include screen readers, braille keyboards and enlargement tools to assist with visual impairments which all require care to be spent labelling design elements, making different elements visually distinct and ensuring layouts are optimal however I wanted to further research the more nuanced strategies required to make my site open to all.

Initially I found this resource: \cite{UnderstandingWCAG222024} - a fantastic series of articles from the UK government explaining not just the legally mandated policies but also how to effectively utilise them for an accessible internet. The provide a number of user 'profiles' to help developers screen their software and see how useable it is to each. This resource focus on the WCAG 2.2 (Web content accessibility guidelines) standards which dictate accessibility requirements all public services must adhere to.

Whilst my artifact would be exempt from this legislation as an academic research product I want to still strive to meet and achieve these targets as that is most beneficial from my users.

There are 4 main principles behind this guidance.

1)
To meet WCAG 2.2 Principle 1: Perceivable you need to make sure users can recognise and use your service with the senses that are available to them.
This means you need to do things like:
\begin{itemize}
    \item provide text alternatives 'alt text' for non-text content
    \item provide transcripts for audio and video
    \item provide captions for video
    \item make sure content is structured logically and can be navigated and read by a screen reader - this also helps if stylesheets are disabled
    \item use the proper markup for every feature (for example, forms and data tables), so the relationships between content are defined properly
    \item not use colour as the only way to explain or distinguish something
    \item use text colours that show up clearly against the background colour
    \item make sure every feature can be used when text size is increased by 200% and that content reflows to a single column when it's increased by 400%
    \item not use images of text
    \item make sure your service is responsive - for example to the user's device, page orientation and font size they like to use
    \item make sure your service works well with assistive technologies - for example, important messages are marked up in a way that the screen readers knows they're important
\end{itemize}
I will meet most of these points purely through developing a fast and performant website as keeping assets as pure text (rather than rendered images) is faster to process for a backend server. Colour is an interesting point mentioned and how it cannot be the only feature to distinguish icons as this contradicts some minimalist design principles (which dictate a single change in appearance should distinguish features), although of course this principles are more important and I will rethink my approach as I build my UI.

2)
To meet WCAG 2.2 Principle 2: Operable, you have to make sure users can find and use your content, regardless of how they choose to access it (for example, using a keyboard or voice commands).
This means you need to do things like:
\begin{itemize}
    \item make sure everything works for keyboard-only users
    \item let people play, pause and stop any moving content
    \item not use blinking or flashing content - or let the user disable animations
    \item provide a 'skip to content' link or equivalent
    \item use descriptive titles for pages and frames
    \item make sure users can move through content in a way that makes sense
    \item use descriptive links so users know where a link will take them, or what downloadable linked content is
    \item use meaningful headings and labels, making sure that any accessible labels match or closely resemble the label you're using in the interface
    \item make it easy for keyboard users to see the item their keyboard or assistive technology is currently focused on - this is known as ‘active focus'
    \item only use things like mouse events or dynamic interactions (like swiping or dragging) when they're strictly necessary - or let the user disable them and interact with the interface in a different way
    \item make it easy for users to disable and change shortcut keys
    \item make sure interactive elements such as buttons are big enough or spaced far enough apart to make it easy to select the right one
\end{itemize}
Once again all of these points are met by my intention to produce a website that is easy to navigate and use by any user, the only potential change to make would be ensuring that throughout my more intricate methods of tracking user progression through a graph, more primitive 'breadcrumbs' are left behind for screen readers etc to follow

3)
To meet WCAG 2.2 Principle 3: Understandable, you have to make sure people can understand your content and how the service works.
This means you need to do things like:
\begin{itemize}
\item make it clear what language the content is written in, and indicate if this changes
\item make sure features look consistent and behave in predictable ways, including help mechanisms
\item make sure all form fields have visible and meaningful labels - and that they're marked up properly
\item make it easy for people to identify and correct errors in forms - you can find best practice for form design in the GOV.UK Design System
\item make it easy for people to re-enter information they've previously entered into a form
\item make it easy for people to log in without having to remember information or solve a problem
\end{itemize}
This third principle has received some criticism from experts due mostly to the last few points. Internet security requires checking that a user is who they say they are through a variety of methods such as password re-entry. Whilst this is frustrating to experience (even more so once captchas etc are added) they are unfortunately necessary to protect user data, even if that doesn't favour accessibility. Of course my authentication forms will still meet standards outlined however making it easy to input details 

4)
To meet WCAG 2.2 Principle 4: Robust, you must make sure your content can be interpreted reliably by a wide variety of user agents (including assistive technologies).

This means you need to do things like:
\begin{itemize}
\item make sure your code lets assistive technologies know what every user interface component is for, what state it's currently in and if it changes
\item make sure important status messages or modal dialogs are marked up in a way that informs users of their presence and purpose, and lets them interact with them using their assistive technology
\end{itemize}

This is the most challenging of the 4 principles as it requires specialised, additional programming to implement. Unfortunately I don't have easy access to external screen reading tools to test my program thoroughly although I can use the dictation tools built into platforms such as Windows and IOS as rudimentary insight. From those learnings I can tweak identifying tags to help make my artifact accessible to all.
